\documentclass[10pt]{article}
\usepackage{multicol}
\usepackage{xltxtra}
\usepackage{marathi}
\दुसराटंक{\mukta}{Mukta}
\दुसराटंक{\modak}{Modak}
\दुसराटंक{\jaini}{Jaini-Regular}
\दुसराटंक{\jainip}{JainiPurva}
\दुसराटंक{\baloo}{Baloo2-Regular}
\दुसराटंक{\gotu}{Gotu}
\usepackage{mdframed}
\usepackage{framed}
\usepackage{tcolorbox}
\usepackage{xcolor}
\usepackage{hyperref}
\hypersetup{
    colorlinks,
    linkcolor={red!50!black},
    citecolor={blue!50!black},
    urlcolor={blue!80!black}
}
\usepackage{capt-of}
\renewcommand{\tablename}{कोष्टक}
\usepackage{standalone}
\makeatletter
\RequirePackage{devanagaridigits}
\def\@arabic#1{\expandafter\devanagaridigits\expandafter{\number#1}}
\makeatother
\usepackage{minted}
\usemintedstyle{bw}
\setmonofont[Script=Devanagari]{Mukta-Regular}
\usepackage{fontawesome5}
\renewcommand{\abstractname}{सारांश}
\renewcommand{\contentsname}{अनुक्रमणिका}
\usepackage{longtable}
\usepackage{setspace}
\setlength{\parindent}{0pt}
\makeatletter
\renewcommand{\@dotsep}{10000}
\makeatother
\title{एक-टाईप टंक}
\author{निरंजन}
 \date{
    \begin{multicols}{2}
        आवृत्ती ०.२ - २६ जुलै, २०२०
  
        \columnbreak
        संस्करण ०.२ - २६ जुलै, २०२०
    \end{multicols}
{\small\faIcon{gitlab}\quad\url{https://gitlab.com/niranjanvikastambe/ektype-tanka}}
}
\sloppy
\hyphenpenalty=10000
\exhyphenpenalty=1000
\setlength\columnsep{20pt}

\begin{document}
 \maketitle
 \अंतरबदल{1.5}
 \begin{abstract}
 \begin{multicols}{2}
 \noindent एक-टाईप ह्या संस्थेच्या काही उत्कृष्ट देवनागरी टंकांचा वापर लाटेक्-सह सहजतेने करता यावा ह्यासाठी ह्या आज्ञासंचाद्वारे मुक्त, जैनी, जैनी पूर्वा, बलू मोदक व गोटू हे टंक पुरवण्यात येत आहेत. हा आज्ञासंच असल्यास \verb|fontspec| चा वापर करून वरीलपैकी कोणताही टंक लाटेक्-मध्ये वापरता येईल. पुढे ह्या सर्व टंकांचे नमुने पुरवण्यात आले आहेत.

 \columnbreak
 एक-टाईप संस्था के कई उत्कृष्ट देवनागरी टंकों का प्रयोग लाटेक् में सुलभ हो इस लिए यह आज्ञासमुच्चय मुक्त, जैनी, जैनी पूर्वा, बलू मोदक और गोटू यह टंक उपलब्ध करवा रहा है। इस आज्ञासमुच्चय के होने पर \verb|fontspec| का प्रयोग कर उपरोक्त कोई भी टंक लाटेक् में प्रयोग किया जा सकता है। इन टंकों के कई उदाहरण इस दस्तावेज़ में पाए जा सकते हैं।
 \end{multicols}
 \end{abstract}
 \tableofcontents\clearpage\pagebreak
 \section{मुक्त}
 \begin{multicols}{2}
 मुक्त हा समानांतर टंक आहे. आज्ञावल्यांच्या लेखनाकरिता हे टंक वापरण्याची पद्धत आहे. अशा वेळी लॅटिन व देवनागरीसाठी दोन टंक आलटून पालटून वापरावे लागतात, परंतु मुक्त हा टंक वापरल्यास ही अडचण येत नाही. पुढील उदाहरणात देवनागरी व लॅटिन मजकूर एकत्र आहे.

 \columnbreak
 मुक्त यह एक समानान्तर टंक आहे। आज्ञावलियों का लेखन करते समय इन टंकों का प्रयोग करने का प्रघात है। तब लैटिन और देवनागरी के लिए दो अलग अलग टंक चुनने पड़ते हैं। लेकिन मुक्त टंक चुनने से यह समस्या नहीं आती। यह उदाहरण देखिए। इस में देवनागरी तथा लैटिन पाठ भी है।
 \end{multicols}

 \begin{framed}
 \begin{minted}[linenos,escapeinside=//]{latex}
\documentclass{article}
\usepackage{marathi}

\begin{document}
\बदल{5} /\%/ ह्या आज्ञेनंतर लिहिलेल्या मजकुरात मूळ अंतराच्या पाचपट अंतर सुटते.
\परिच्छेद
\end{document}
 \end{minted}
 \end{framed}
 \begin{multicols}{2}
 वेगवेगळ्या आकारांत मुक्त असा दिसतो.\\

 \columnbreak
 भिन्न आकारों में मुक्त इस तरह दिखता है।\\
 \end{multicols}
 {\centering\mukta
 \tiny नमस्कार\\    
 \scriptsize नमस्कार\\
 \footnotesize नमस्कार\\
 \small नमस्कार\\
 \large नमस्कार\\
 \Large नमस्कार\\
 \LARGE नमस्कार\\
 \Huge नमस्कार\\}
 \section{बलू}
 \begin{multicols}{2}
 बलू हा एक समरेखा टंक आहे. निरनिराळ्या आकारांत तो असा दिसतो.

 \columnbreak
 बलू एक समरेखा टंक है। भिन्न आकारों में बलू कुछ इस तरह दिखता है।\\
 \end{multicols}
 {\centering\baloo
 \tiny नमस्कार\\    
 \scriptsize नमस्कार\\
 \footnotesize नमस्कार\\
 \small नमस्कार\\
 \large नमस्कार\\
 \Large नमस्कार\\
 \LARGE नमस्कार\\
 \Huge नमस्कार\\}
 \section{जैनी}
 \begin{multicols}{2}
 जैन कल्पसूत्रांच्या शैलीवर आधारित हा टंक विविध आकारांमध्ये असा दिसतो

 \columnbreak
 जैन कल्पसूत्रों के शैली पर आधारित यह टंक भिन्न आकारों में ऐसा दिखाई देता है।\\
 \end{multicols}
 {\centering\jaini
 \tiny नमस्कार\\    
 \scriptsize नमस्कार\\
 \footnotesize नमस्कार\\
 \small नमस्कार\\
 \large नमस्कार\\
 \Large नमस्कार\\
 \LARGE नमस्कार\\
 \Huge नमस्कार\\}
 \subsection{जैनी पूर्वा}\vspace*{-0.2cm}
 \begin{multicols}{2}
 हा टंक कल्पसूत्रांतील शैलीचे हुबेहुब प्रतिरूप आहे, परंतु सर्वसामान्य वाचकास वाचण्यासाठी किंचित अवघड आहे. पुरातन ग्रंथांच्या लेखनाकरिता हा टंक खूपच उपयुक्त आहे.

 \columnbreak
 यह टंक कल्पसूत्रों के अक्षर का साक्षात प्रतिरूप है, किन्तु सर्वसामान्य वाचक को पढ़ने के लिए थोड़ा मुश्किल है। पुरातन ग्रंथों के लेखन के लिए यह टंक बहुत ही उपयुक्त है।
 \end{multicols}
 \section{मोदक}
 \begin{multicols}{2}
 हा दर्शनी टंक शीर्षकांकरिता वापरला जातो. ह्याचे एकच वजन उपलब्ध आहे. सर्व आकारांमध्ये हा टंक पुढीलप्रमाणे दिसतो.

 \columnbreak
 इस दर्शनी टंक का प्रयोग शीर्षकों के लिए होता है। यह टंक एक ही वज़न में उपलब्ध है। इस टंक के सर्व आकार इस तरह दिखते हैं।
 \end{multicols}
 {\centering\modak
 \tiny नमस्कार\\
 \scriptsize नमस्कार\\
 \footnotesize नमस्कार\\
 \small नमस्कार\\
 \large नमस्कार\\
 \Large नमस्कार\\
 \LARGE नमस्कार\\
 \Huge नमस्कार\\}
 \pagebreak
 \section{गोटू}
 \begin{multicols}{2}
 हा एक पारंपरिक टंक आहे.\linebreak ह्याच्या अक्षरांची लांबी कमी जास्त होते.
 
 \columnbreak
 यह एक पारम्परिक टंक है। इस \linebreak के अक्षरों की चौड़ाई अस्थिर होती है।
 \end{multicols}
 \centering\gotu
 \tiny नमस्कार\\
 \scriptsize नमस्कार\\
 \footnotesize नमस्कार\\
 \small नमस्कार\\
 \large नमस्कार\\
 \Large नमस्कार\\
 \LARGE नमस्कार\\
 \Huge नमस्कार\\
\end{document}